\chapter{Thuật ngữ}
\label{app:glossary}

\newcommand{\gloss}[2]{\item[#1] #2}
\begin{description}\setlength{\itemsep}{2pt}
\gloss{Actuator}{Các endpoint vận hành của Spring Boot
  (\code{/actuator/health}, \code{/metrics}, \code{/prometheus}); mở ra
  và phân quyền là hai quyết định riêng (Ch.~1).}
\gloss{Advertised listener}{Địa chỉ mà broker Kafka bảo client dùng
  sau bước bootstrap --- nguyên nhân thường gặp của hiện tượng kết nối
  được rồi timeout (Ch.~10).}
\gloss{Algorithm confusion}{(nhầm lẫn thuật toán) Tấn công sửa header
  \code{alg} của token để bên xác minh áp dụng HMAC với khóa công khai
  làm secret; bị chặn khi loại khóa, không phải header, chọn thuật toán
  (Ch.~28).}
\gloss{Artifact}{Bất cứ thứ gì một build sinh ra đáng giữ lại (jar,
  image, báo cáo); phải bất biến và truy nguyên được (Ch.~21).}
\gloss{Auto-configuration}{(tự cấu hình) Spring Boot áp dụng cấu hình
  có điều kiện dựa trên classpath, thuộc tính và bean của người dùng
  (Ch.~1).}
\gloss{Backend for frontend}{Service nhỏ ghép nhiều lời gọi service
  thành một response theo hình dạng client; chỗ đúng cho việc gộp mà
  gateway không được làm (Ch.~31).}
\gloss{Bootstrap servers}{Các địa chỉ Kafka ban đầu mà client liên hệ
  để lấy metadata; mọi lưu lượng sau đó dùng advertised listener
  (Ch.~10).}
\gloss{Bulkhead}{Giới hạn số lời gọi đồng thời tới một phụ thuộc
  (semaphore hoặc thread pool riêng) để một phụ thuộc chậm không tiêu
  hết thread của bên gọi (Ch.~32).}
\gloss{Cipher value}{(giá trị mã hóa) Giá trị config-server có tiền tố
  \code{\{cipher\}}, mã hóa bằng khóa của server để repository chỉ giữ
  bản mã; được giải mã lúc phục vụ (Ch.~29).}
\gloss{Circuit breaker}{(bộ ngắt mạch) Máy trạng thái trên cửa sổ trượt
  các kết quả gọi: CLOSED cho qua, OPEN từ chối tức thì trong một
  khoảng chờ, HALF\_OPEN thăm dò bằng vài lời gọi (Ch.~32).}
\gloss{ClusterIP}{Kiểu Service mặc định: một IP ảo ổn định, chỉ tới
  được từ bên trong cụm (Ch.~18).}
\gloss{Consumer group}{(consumer group) Một nhóm consumer Kafka có tên,
  chia nhau các partition của một topic; mỗi nhóm giữ offset riêng
  (Ch.~10).}
\gloss{Consumer lag}{(độ trễ tiêu thụ) Khoảng cách giữa offset mới nhất
  của partition và offset đã commit của group; đo bằng giây, đặt cảnh
  báo trên nó (Ch.~34).}
\gloss{Controller (k8s)}{Một vòng lặp điều hòa hội tụ trạng thái thực
  tế về trạng thái mong muốn cho một kind object (Ch.~14).}
\gloss{CORS / preflight}{Cơ chế trình duyệt cưỡng chế để nới lỏng
  chính sách same-origin; request cross-origin không đơn giản tốn
  thêm một vòng \code{OPTIONS} trừ khi được cache --- và không bao giờ
  xảy ra khi trang và API chung một origin (Ch.~31, 38).}
\gloss{CQRS}{Tách lệnh/truy vấn xuyên service: ghi đi tới service sở
  hữu, đọc đi tới một mô hình dựng từ sự kiện của nó (Ch.~33).}
\gloss{CrashLoopBackOff}{Trạng thái pod: container liên tục thoát; độ
  trễ khởi động lại tăng gấp đôi mỗi lần (Ch.~24).}
\gloss{CRD / Operator}{Các kind resource tùy biến cộng một controller
  cho chúng --- cách Strimzi quản lý Kafka theo kiểu khai báo
  (Ch.~25).}
\gloss{Dead-letter topic}{(dead-letter topic) Nơi consumer chuyển những
  thông điệp xử lý thất bại nhiều lần (Ch.~11).}
\gloss{Deployment}{Controller workload cho các replica thay thế được
  cho nhau; quản lý ReplicaSet; cho phép cập nhật cuốn chiếu và
  rollback (Ch.~15).}
\gloss{Digest}{Mã băm nội dung bất biến của một image
  (\code{@sha256:...}); tag dịch chuyển, digest thì không (Ch.~13).}
\gloss{Dual write}{(ghi kép) Ghi vào hai hệ thống (cơ sở dữ liệu và
  Kafka) không có transaction chung; cái này có thể thành công trong khi
  cái kia thất bại (Ch.~33).}
\gloss{Endpoint (k8s)}{Một pod Ready đứng sau một Service; danh sách
  sống được duy trì theo label selector và readiness (Ch.~18).}
\gloss{Eureka}{Service registry của Netflix/Spring Cloud: đăng ký,
  heartbeat, cache phía client (Ch.~3).}
\gloss{Event loop}{(vòng lặp sự kiện) Nhóm thread nhỏ, cố định của
  Reactor Netty, mỗi thread ghép kênh nhiều socket bằng I/O không
  chặn; một lời gọi chặn làm khựng mọi request trên vòng lặp đó
  (Ch.~31).}
\gloss{explain (Mongo)}{\code{find(...).explain("executionStats")}:
  kế hoạch truy vấn với các stage \code{IXSCAN}/\code{FETCH} và
  \code{totalDocsExamined} so với \code{nReturned} --- thứ đầu tiên
  cần đọc khi một truy vấn chậm (Ch.~36).}
\gloss{Fallback}{Giá trị hay hành động thay thế khi lời gọi được bảo vệ
  thất bại hoặc bị từ chối; nên là giá trị cache, suy giảm-nhưng-có
  kiểu, rỗng, hoặc lỗi có kiểu --- không bao giờ \code{null} (Ch.~32).}
\gloss{Flyway}{Các migration SQL có phiên bản, bất biến, áp dụng lúc
  khởi động và ghi vào một bảng lịch sử (Ch.~7).}
\gloss{GitOps}{Mô hình triển khai trong đó một agent trong cụm kéo
  trạng thái mong muốn từ git và hội tụ cụm về đó (Ch.~26).}
\gloss{Head-based / tail-based sampling}{(lấy mẫu theo đầu / theo
  đuôi) Quyết định tại span gốc có giữ trace hay không (rẻ, mù với kết
  quả) so với quyết định sau khi biết trọn trace (giữ lỗi và request
  chậm; cần collector) (Ch.~37).}
\gloss{HS256 / RS256}{Thuật toán ký JWT. HS256: một secret dùng chung
  vừa ký vừa xác minh. RS256: khóa riêng RSA ký, khóa công khai xác
  minh --- một bên phát hành, bao nhiêu bên xác minh cũng được
  (Ch.~28).}
\gloss{HttpOnly cookie}{Cookie mà JavaScript không đọc được; cùng
  \code{Secure}, \code{SameSite} và một \code{Path} hẹp, nó là chỗ ở
  an toàn cho refresh token (Ch.~38).}
\gloss{Idempotent consumer}{(consumer lũy đẳng) Listener giành mỗi id
  sự kiện trong cùng transaction với hiệu ứng phụ, nên giao lại là
  no-op (Ch.~34).}
\gloss{Idempotent producer}{(producer lũy đẳng) Producer Kafka mà các
  lần thử lại được broker khử trùng theo số thứ tự; mặc định từ
  Kafka~3.0 (Ch.~34).}
\gloss{Ingress}{Object định tuyến HTTP tầng 7 (host/đường dẫn
  $\rightarrow$ Service), được thực thi bởi một Ingress controller như
  Traefik (Ch.~18).}
\gloss{Jib}{Plugin Maven build image OCI theo layer mà không cần Docker
  daemon (Ch.~13).}
\gloss{JWKS}{JSON Web Key Set: các khóa công khai mà bên phát hành công
  bố tại một URL quy ước để bên xác minh tải và cache (Ch.~28).}
\gloss{JWT}{Token danh tính tự chứa, có chữ ký:
  header.payload.signature; kiểm tra được mà không cần cơ sở dữ liệu
  (Ch.~5; ở quy mô lớn, Ch.~28).}
\gloss{Keyset pagination}{(phân trang theo khóa) Phân trang bằng ``các
  dòng sau khóa này'' trên cột sort có index thay vì \code{skip}, để
  trang sâu tốn bằng trang đầu (Ch.~36).}
\gloss{kid}{Key id trong header JWT, đặt tên khóa đã ký nó; thứ cho
  phép xoay khóa trong khi token cũ còn lưu hành (Ch.~28).}
\gloss{KRaft}{Cơ chế đồng thuận Raft tích hợp sẵn của Kafka thay cho
  ZooKeeper (Ch.~10).}
\gloss{kubelet}{Agent trên mỗi node, chạy pod và thực thi probe
  (Ch.~14).}
\gloss{Kustomize}{Ghép manifest không dùng template: base, overlay,
  patch; tích hợp sẵn trong kubectl (Ch.~20).}
\gloss{Lease (Eureka)}{Đăng ký của một instance, gia hạn bằng heartbeat
  mỗi 30\,s và hết hạn sau 90\,s im lặng; server quét mỗi 60\,s, nên
  instance chết có thể lưu lại vài phút (Ch.~30).}
\gloss{Liveness probe}{``Khởi động lại tôi nếu cái này hỏng'' --- probe
  cần dè dặt nhất (Ch.~17).}
\gloss{Multipart upload}{Cơ chế S3 chia một object lớn thành các phần
  ít nhất 5\,MB, upload song song và ráp phía server; phần bị bỏ dở cần
  quy tắc lifecycle hủy (Ch.~35).}
\gloss{Namespace}{Phân vùng của cụm; phạm vi cho tên và RBAC
  (Ch.~14).}
\gloss{OOMKilled}{Exit 137: kernel giết container vì vượt giới hạn bộ
  nhớ; không có trace nào từ phía ứng dụng (Ch.~17).}
\gloss{Outbox pattern}{Làm cho ghi DB + phát sự kiện trở nên nguyên tử
  trên thực tế bằng cách ghi sự kiện vào một bảng trong cùng giao dịch
  (Ch.~11; Ch.~33).}
\gloss{Partition}{(phân vùng) Đơn vị thứ tự và song song của Kafka;
  một chuỗi bản ghi bất biến, có thứ tự (Ch.~10).}
\gloss{Path-style addressing}{(địa chỉ kiểu đường dẫn) Dạng URL S3 với
  bucket nằm trong đường dẫn (\code{host/bucket/key}) thay vì là
  subdomain DNS; MinIO cần dạng này khi không có wildcard DNS
  (Ch.~35).}
\gloss{Phantom token}{Định danh mờ bên ngoài, JWT đầy đủ bên trong;
  gateway đổi cái này lấy cái kia qua một kho nhanh, cho token ngoài
  mỏng và thu hồi tức thì với cái giá là một phụ thuộc mới (Ch.~28).}
\gloss{Poison pill}{Một thông điệp giải tuần tự thất bại mãi mãi, làm
  kẹt partition của nó trừ khi xử lý lỗi bỏ qua nó (Ch.~11).}
\gloss{Presigned URL}{(presigned URL) Một link có chữ ký, giới hạn thời
  gian, cho phép trình duyệt truy cập trực tiếp một object trên
  S3/MinIO (Ch.~9).}
\gloss{Probe (startup/readiness/liveness)}{Các kiểm tra sức khỏe của
  kubelet, lần lượt gác việc khởi động lại, nhận lưu lượng, và thời
  gian ân hạn lúc boot (Ch.~17).}
\gloss{Property source precedence}{(thứ tự ưu tiên nguồn thuộc tính)
  Danh sách có thứ tự các nguồn cấu hình của Spring Boot; biến môi
  trường đứng trên tài liệu từ config-server, tài liệu đó đứng trên
  YAML đóng gói (Ch.~29).}
\gloss{PV / PVC / StorageClass}{Lưu trữ đã cấp phát, yêu cầu của một
  workload lên nó, và nhà máy tạo PV theo nhu cầu (Ch.~19).}
\gloss{RBAC}{Kiểm soát truy cập theo vai trò: ServiceAccount (danh
  tính) + Role (động từ trên resource) + RoleBinding (mối nối)
  (Ch.~22).}
\gloss{Reconciliation}{(điều hòa) Vòng lặp điều khiển: so sánh trạng
  thái mong muốn và thực tế, hành động trên phần chênh lệch, mãi mãi
  (Ch.~14).}
\gloss{RED / USE}{Rate, Errors, Duration cho service theo request;
  Utilization, Saturation, Errors cho tài nguyên --- hai bộ từ vựng
  dashboard (Ch.~37).}
\gloss{Refresh scope}{Các bean được dựng lại khi \code{/actuator/refresh}
  để thay đổi cấu hình có hiệu lực không cần khởi động lại; theo từng
  instance trừ khi Spring Cloud Bus phát sóng (Ch.~29).}
\gloss{Registry}{Kho lưu image nói giao thức Docker/OCI; giao diện giữa
  bên build và cụm (Ch.~13).}
\gloss{Relaxed binding}{Cách Spring ánh xạ
  \code{SPRING\_DATASOURCE\_URL} thành \code{spring.datasource.url}
  (Ch.~1).}
\gloss{ReplicaSet}{Bộ đếm giữ N pod của một template luôn sống; do
  Deployment quản lý, giữ lại để rollback (Ch.~15).}
\gloss{Retry storm}{(bão retry) Các lần retry đồng bộ từ nhiều client
  nhân tải lên một phụ thuộc đang lỗi; ngăn bằng jitter, ngân sách
  retry và chỉ retry lời gọi lũy đẳng (Ch.~32).}
\gloss{Retry topic}{Topic theo độ trễ mà một record thất bại được phát
  lại vào (\code{-retry-2000}, \ldots), kết thúc ở dead-letter topic;
  \code{@RetryableTopic} của Spring (Ch.~34).}
\gloss{Rolling update}{(cập nhật cuốn chiếu) Thay pod từng phần, gác
  bằng readiness, để service không bao giờ dừng hẳn (Ch.~15).}
\gloss{Saga}{Một lần ghi xuyên service dưới dạng chuỗi transaction cục
  bộ, mỗi bước có hành động bù nếu bước sau thất bại (Ch.~33).}
\gloss{Secret (k8s)}{Object key/value theo namespace cho thông tin xác
  thực; mã hóa base64, được bảo vệ bởi RBAC chứ không phải mật mã học
  (Ch.~16).}
\gloss{Self-preservation}{(chế độ tự bảo toàn) Chế độ Eureka vào khi
  gia hạn rơi dưới 85\,\% kỳ vọng; loại bỏ dừng hoàn toàn --- kích hoạt
  chỉ bởi một crash trên đội tàu năm instance (Ch.~30).}
\gloss{Service (k8s)}{Tên ổn định + IP ảo trên một tập pod chọn theo
  label, lọc theo readiness (Ch.~18).}
\gloss{ServiceAccount}{Một danh tính phi con người để truy cập API ---
  thứ Jenkins dùng để deploy (Ch.~22).}
\gloss{SigV4}{AWS Signature Version 4: một chuỗi HMAC trên ngày, region
  và dịch vụ ký một canonical request luôn bao gồm header \code{Host}
  --- presigned URL bị gắn chặt với host nó được ký cho (Ch.~35).}
\gloss{Single-flight refresh}{(refresh một chuyến) Một promise chờ
  dùng chung để các 401 đồng thời chỉ kích hoạt đúng một lần refresh
  token và tất cả replay với kết quả của nó; bắt buộc khi refresh
  token có xoay vòng (Ch.~38).}
\gloss{StatefulSet}{Controller workload cho các replica có trạng thái:
  tên ổn định, PVC riêng từng replica, thay thế theo thứ tự (Ch.~15).}
\gloss{Text index (Mongo)}{Index tách token và rút gốc các trường chuỗi
  có trọng số, truy vấn qua \code{\$text}/\code{TextCriteria}; khớp từ
  trọn vẹn, không khớp chuỗi con (Ch.~36).}
\gloss{Token bucket}{Thuật toán rate limit: token nạp lại đều đặn tới
  một sức chứa burst và mỗi request tiêu một; Spring Cloud Gateway giữ
  bucket trong Redis (Ch.~31).}
\gloss{TOTP}{Mật khẩu một lần theo thời gian: bí mật chia sẻ + đồng hồ,
  RFC 6238 (Ch.~5).}
\gloss{traceparent}{Header W3C Trace Context mang trace id, span id
  cha và cờ lấy mẫu giữa các service; Boot~3 phát nó và cũng chấp nhận
  B3 (Ch.~37).}
\gloss{Trusted edge}{(biên tin cậy) Gateway kiểm tra JWT một lần và
  chuyển tiếp các header danh tính; các service tin chúng --- chỉ đúng
  khi không có đường nào đi vòng qua gateway (Ch.~5).}
\end{description}
